\section{Introduction}

Recent advances in educational AI have expanded both the assessment of learner work and the modeling of individual learning behavior. Multimodal models support educator-aligned assessment of artistic performance \citep{ye2026multimodal}, while history-conditioned models infer student knowledge and behavioral tendencies from previous interactions \citep{duan2026history}. Individualized simulators further model learner responses and their adjustment under instructional guidance \citep{yang2026simulator}. Translating these capabilities into personalized educational support requires establishing which inferred characteristics are relevant to instruction. \textbf{Accurate descriptions of individual differences do not, on their own, determine appropriate instructional priorities.} \textbf{The central challenge is therefore to identify pedagogically relevant individual characteristics from limited learning histories and use them to guide individualized instructional support.}

Recent methods learn individual representations primarily through response prediction or preference matching. Student profiles are optimized to predict subsequent dialogue turns \citep{duan2026history}, while learned user summaries are optimized for personalized preference prediction \citep{nam2026summaries}. \textbf{Although these objectives identify useful personal information, they do not explicitly supervise how a supported learner characteristic should influence an instructional decision.} Simulator-based tutor optimization advances the instructional application, but the experiment reported by \citet{yang2026simulator} uses a pooled simulator rather than individually specialized models, leaving the contribution of personal histories to improved guidance untested in that setting. A more targeted learning objective must account for both the basis of a learner-specific inference and its instructional relevance. \textbf{Process records document prior behavior, whereas expert judgments evaluate the appropriateness of support. Connecting these judgments to the observations that justify them provides a basis for learning which individual characteristics are supported and how they should inform instruction.}

\textbf{To investigate this problem, we develop a digital creative-learning platform and design a multimodal dataset connecting individual creation histories with instructional support.} The platform supports drawing activities with a defined feedback checkpoint and an optional revision stage, preserving intermediate artworks and editing sequences before and after assistance. Activity design draws on Studio Thinking's account of artistic habits, including envisioning, exploration, and reflection \citep{hetland2013studio}, and on expert-developed activities associated with the National Core Arts Standards and its Model Cornerstone Assessments.\footnote{\url{https://www.nationalartsstandards.org}} Tasks are adapted to participants' educational levels and artistic experience, with expert review and piloting examining construct validity and developmental appropriateness. \textbf{Educators assess instructional responses in relation to each learner's recorded work and identify the observations supporting their judgments.} The dataset thus connects demonstrated learning behavior with expert assessments of appropriate assistance.

\textbf{Using this resource, we investigate learner representation learning for personalized instructional support.} The aim is to capture individual capabilities and strategy patterns that are both grounded in recorded experience and relevant to subsequent assistance. Process-based modeling and expert pedagogical judgments provide complementary information for studying how these characteristics can be identified and used. Feedback serves as an initial application through which to examine their instructional value, without restricting the representation to a single downstream task. \textbf{Evaluation focuses on whether inferred characteristics are supported by the available history and whether they contribute to more appropriate learner-specific support.}
